% ============================================================
% Section 24: 一维惰性气体链的 Lennard-Jones 势与线弹性模量
% ============================================================
\section{固体理论：一维惰性气体链——Lennard-Jones 势、晶格能与线弹性模量}
\label{sec:lj-chain}

\begin{modelbox}
$N$ 个惰性气体原子组成线性布拉维固体链（一维单原子链），平均每两原子间相互作用势为
\[
  \varphi(x) = \varphi_0 \left[ \left(\frac{\sigma}{x}\right)^{12} - 2\left(\frac{\sigma}{x}\right)^{6} \right].
\]
试求：
\begin{enumerate}[nosep]
  \item 原子间的平均距离 $x_0$；
  \item 每个原子的平均晶格能 $u_0$；
  \item 线弹性模量 $K$。
\end{enumerate}
\end{modelbox}

\vspace{0.3em}
\begin{insightbox}{结论预告}
\[
  x_0 = \sigma, \qquad u_0 = -\varphi_0, \qquad K = \frac{72\varphi_0}{\sigma}.
\]
本题是典型的\textbf{求极小值 $\to$ 读函数值 $\to$ 求曲率}三步曲。势函数的参数化经过特殊设计（$-2$ 而非 $-1$），使得平衡位置恰好等于 $\sigma$，最小势能恰好等于 $-\varphi_0$。
\end{insightbox}

\subsection{关键观察：势函数的配方}

令 $y = (\sigma/x)^6$，则
\[
  \varphi(x) = \varphi_0\,(y^2 - 2y) = \varphi_0\,[(y-1)^2 - 1].
\]
这是关于 $y$ 的抛物线，\textbf{极小值显然在 $y=1$ 处}。这种改写让求根变成口算，避免繁琐的幂次运算。

\subsection{平衡距离 $x_0$}

物理条件：平衡时原子所受合力为零，即势能取极小值
\[
  \left.\frac{d\varphi}{dx}\right|_{x_0} = 0.
\]

求导：
\begin{align*}
  \frac{d\varphi}{dx}
  &= \varphi_0 \left[ -12\frac{\sigma^{12}}{x^{13}} + 12\frac{\sigma^6}{x^7} \right] \\
  &= \frac{12\varphi_0}{x} \left[ \left(\frac{\sigma}{x}\right)^6 - \left(\frac{\sigma}{x}\right)^{12} \right].
\end{align*}

令其为零（$x\neq 0$），得
\[
  \left(\frac{\sigma}{x_0}\right)^6 = 1 \quad\Longrightarrow\quad \boxed{x_0 = \sigma}.
\]

验证极小：$\varphi''(\sigma) = 72\varphi_0/\sigma^2 > 0$（见下文），确为极小值。

\begin{tipbox}{与标准 Lennard-Jones 势的区分}
标准 LJ 势常写作
\[
  \varphi_{\text{LJ}}(r) = 4\varepsilon\left[\left(\frac{\sigma}{r}\right)^{12} - \left(\frac{\sigma}{r}\right)^6\right],
\]
其极小值在 $r_0 = 2^{1/6}\sigma$ 处，$\varphi_{\min} = -\varepsilon$。

本题为 $-2(\sigma/x)^6$，相当于把吸引项加倍，使极小值恰好落在 $x_0 = \sigma$ 处。\textbf{考场上务必看清系数}，不要凭肌肉记忆写 $2^{1/6}\sigma$。
\end{tipbox}

\subsubsection*{概念铺垫：键能、晶格能与分摊法则}

\begin{insightbox}{从``手拉手''模型理解键能与晶格能}
想象 $N$ 个人排成一排，每人两只手分别与左右邻居``手拉手''。
\begin{itemize}[nosep]
  \item \textbf{键能}：任意相邻两人之间``手拉手纽带''的能量。拉得太近会排斥（能量升高），拉得太远会吸引（能量也升高），在某个舒服距离上能量最低。这个最低能量就是一对原子的\textbf{键能} $\varphi(x_0)$。
  \item \textbf{晶格能}：整条人链中\emph{所有}``手拉手纽带''能量的总和。注意：每对邻居只算\emph{一次}，不能重复计数。
\end{itemize}
\end{insightbox}

\begin{modelbox}{两种等价的计数方法}
\textbf{方法一：数键法（推荐）}

$N$ 个原子排成一维链，最近邻对（键）的总数为 $N-1 \approx N$（忽略边界）。
\[
  U_{\text{总}} = (N-1)\,\varphi(x_0) \approx N\varphi(x_0), \qquad u_0 = \frac{U_{\text{总}}}{N} = \varphi(x_0).
\]

\textbf{方法二：数邻居法（通用公式）}

每个原子有 $z$ 个最近邻（配位数）。一维链 $z=2$，三维 fcc $z=12$。
每个原子``觉得''自己有 $z$ 条键，但\textbf{每条键被 2 个原子共享}，因此实际分摊为
\[
  u_0 = \frac{1}{2} \times z \times \varphi(x_0) = \frac{z}{2}\,\varphi(x_0).
\]

两种方法等价。一维链 $z=2$，恰好 $z/2=1$，所以 $u_0 = \varphi(x_0)$。
\end{modelbox}

\subsection{每个原子的平均晶格能 $u_0$}

\begin{modelbox}[计数规则：一维链的键能分摊]
\begin{itemize}[nosep]
  \item 每个原子有 $\bm{z=2}$ 个最近邻（左、右各 1 个）；
  \item 每对近邻相互作用被 \textbf{2 个原子共享}。
\end{itemize}
因此每个原子平均分摊到的键能为
\[
  u_0 = \frac{1}{2} \times z \times \varphi(x_0) = \frac{z}{2}\,\varphi(x_0).
\]
\end{modelbox}

对一维单原子链，$z=2$，故 $u_0 = \varphi(x_0)$。代入 $x_0 = \sigma$：
\[
  \varphi(\sigma) = \varphi_0(1 - 2) = -\varphi_0 \quad\Longrightarrow\quad \boxed{u_0 = -\varphi_0}.
\]

\begin{tipbox}{常见错误}
\begin{center}
\renewcommand{\arraystretch}{1.4}
\begin{tabular}{|l|l|l|}
\hline
\textbf{错误写法} & \textbf{错误原因} & \textbf{正确理解} \\
\hline
$u_0 = \frac{1}{2}\varphi(x_0)$ & 忘了乘最近邻数 $z=2$ & 每个原子有两个邻居 \\
$u_0 = 2\varphi(x_0)$ & 忘了除以共享因子 2 & 每对键被两个原子分摊 \\
\hline
\end{tabular}
\end{center}
一维链恰好 $z/2 = 1$，结果上看就是 $\varphi(x_0)$ 本身，但这只是巧合。
\end{tipbox}

\subsubsection*{概念铺垫：应变、应力与弹性模量}

\begin{insightbox}{从橡皮筋理解应变与应力}
拿一根原长为 $L_0$ 的橡皮筋：
\begin{itemize}[nosep]
  \item \textbf{应变}（strain）：你把它拉长了 $\Delta L$，\textbf{相对}伸长了多少？
  \[
    \varepsilon = \frac{\Delta L}{L_0} \qquad \text{（无量纲，``百分比''）}
  \]
  应变描述的是\textbf{形变的程度}，与橡皮筋的绝对长度无关。

  \item \textbf{应力}（stress）：橡皮筋内部产生了一个``抵抗''你拉长的恢复力。应力是这个恢复力的\textbf{强度}，即单位面积上的力。
  \[
    \sigma = \frac{F}{A} \qquad \text{（量纲：力/面积 = 压强）}
  \]
\end{itemize}
\end{insightbox}

\begin{modelbox}{胡克定律与弹性模量}
在弹性范围内，应力与应变成正比：
\[
  \sigma = Y \, \varepsilon \quad\Longleftrightarrow\quad \text{弹性模量 } Y = \frac{\sigma}{\varepsilon}.
\]
\begin{itemize}[nosep]
  \item 弹性模量 $Y$ 衡量材料的\textbf{刚度}：$Y$ 越大，同样的应变需要更大的应力（材料越硬）。
  \item 三维中 $Y$ 叫杨氏模量（Young's modulus），量纲为压强（Pa）。
  \item 一维链中没有``横截面积''的概念，\textbf{线弹性模量} $K$ 的定义改为``单位应变所需的恢复力''，量纲为力（N）。
\end{itemize}
\end{modelbox}

\subsection{线弹性模量 $K$}

\subsubsection{定义推导}

一维链的\textbf{线弹性模量}（单位应变所需的恢复力）定义为
\[
  K = x_0 \left.\frac{d^2\varphi}{dx^2}\right|_{x_0}.
\]

\textbf{这个公式怎么来的？} 设链发生均匀应变 $\varepsilon$，键长变为 $x = x_0(1+\varepsilon)$。总晶格能展开：
\[
  U \approx N\varphi(x_0) + N\cdot\frac{1}{2}\varphi''(x_0)\cdot(x_0\varepsilon)^2.
\]
单位长度弹性能：
\[
  \frac{U_{\text{弹}}}{L} = \frac{1}{2}\underbrace{[x_0\,\varphi''(x_0)]}_{K}\,\varepsilon^2.
\]
对比弹性力学 $u_{\text{弹}} = \frac{1}{2}K\varepsilon^2$，即得上式。

\subsubsection{计算二阶导}

\[
  \frac{d^2\varphi}{dx^2} = \varphi_0 \left[ 156\frac{\sigma^{12}}{x^{14}} - 84\frac{\sigma^6}{x^8} \right].
\]

在 $x = \sigma$ 处：
\[
  \left.\frac{d^2\varphi}{dx^2}\right|_{\sigma}
  = \frac{\varphi_0}{\sigma^2}(156 - 84)
  = \frac{72\varphi_0}{\sigma^2}.
\]

于是
\[
  K = \sigma \cdot \frac{72\varphi_0}{\sigma^2} = \boxed{\frac{72\varphi_0}{\sigma}}.
\]

\begin{tipbox}{量纲检验}
$[\varphi_0] = \text{J}$，$[\sigma] = \text{m}$，故 $[K] = \text{J/m} = \text{N}$（力）。

恰为``力/应变''的量纲（应变无量纲）。正确。
\end{tipbox}

\subsection{物理图像与拓展}

\subsubsection{与三维惰性气体晶体的对比}

\begin{center}
\renewcommand{\arraystretch}{1.5}
\begin{tabular}{|l|c|c|}
\hline
\textbf{物理量} & \textbf{一维链（本题）} & \textbf{三维 fcc（参考）} \\
\hline
最近邻数 $z$ & 2 & 12 \\
每个原子晶格能 & $u_0 = \varphi(x_0) = -\varphi_0$ & $u_0 = \frac{z}{2}\varphi(r_0) = -6\varphi_0$（同势） \\
弹性模量 & $K = x_0\varphi''(x_0)$ & 体积模量 $B = \frac{r_0^2}{9V_0}\varphi''(r_0)$ \\
\hline
\end{tabular}
\end{center}

\subsubsection{与晶格振动的联系}

简谐近似下力常数 $\beta = \varphi''(x_0) = 72\varphi_0/\sigma^2$。一维单原子链的声子色散：
\[
  \omega(q) = 2\sqrt{\frac{\beta}{m}}\left|\sin\frac{q\sigma}{2}\right|
            = 2\sqrt{\frac{72\varphi_0}{m\sigma^2}}\left|\sin\frac{q\sigma}{2}\right|.
\]
长波极限 $q\to 0$ 的声速 $c = \sigma\sqrt{\beta/m} = \sqrt{K/m}$。

\subsubsection{Einstein 近似下的振动频率}

若把每个原子看成在平衡位置附近独立振动，则角频率
\[
  \omega_E = \sqrt{\frac{\beta}{m}} = \sqrt{\frac{72\varphi_0}{m\sigma^2}}.
\]
